\noindent We sincerely thank you for your time, effort, and valuable feedback. Your detailed comments and constructive suggestions have significantly improved the paper. We have carefully considered each point and prepared a thorough, point-by-point response. Below, we respond to each comment individually; for ease of reading, the original comments are italicized and displayed in red. Before addressing each comment, we begin with a brief summary of the major revisions.

\begin{itemize}
\item \textbf{Extended sample.} Following R1's and R6's suggestions, we extended our panel data coverage from seven years (1930--1936) to fifteen years (1926--1940). This extension serves two primary purposes:
\begin{enumerate}
    \item Examining additional pre-treatment years (1926--1932 rather than just 1930--1932) enables a more rigorous evaluation of the parallel trends assumption underlying our difference-in-differences design.
    \item The recession of 1937--1938 provides a natural placebo test: if our gold-clause exposure measure captured something other than the mechanism we propose---such as general sensitivity to economic downturns---we would expect it to affect investment during this later period of distress as well.
\end{enumerate}
Figure \ref{fig:rev_gold_coeffs} (Figure 3 in the revised manuscript) plots the year-by-year interaction coefficients from our main investment regression over the extended sample. The figure confirms the absence of pre-trends before 1933, sharp negative effects during the 1933--1934 litigation period, and no effect during the 1937--1938 recession.
\begin{figure}[h]
    \centering
    \includegraphics[width=0.75\textwidth]{../../../../manuscript/figures/body/parallel_trend_plot.pdf}
    \caption{Gold clause effect on investment by year}
    \label{fig:rev_gold_coeffs}
    \footnotesize \textit{Notes.} This figure plots point estimates and 95\% confidence bands from a regression of net investment on $Q$ and $\tilde{d}$ interacted with year indicators, including firm and year fixed effects.
\end{figure}

\item \textbf{Data reconstruction.} While extending the sample, we identified that some scans purchased from Mergent were of poor quality, leading to firm omissions and data entry errors. To address these concerns, we undertook a complete data reconstruction:
\begin{enumerate}
    \item We scanned, in high resolution, the Moody's Industrial Manuals for all sixteen years, ensuring complete legibility of all pages.
    \item We recollected the entire dataset using a double-blind procedure in which two independent research assistants collected data separately, with discrepancies resolved by a third party.
    \item We implemented extensive quality control checks to identify and correct any remaining issues.
\end{enumerate}
This reconstruction process required nearly two years to complete but was essential for ensuring the reliability of our findings.

\item \textbf{Gold clause exposure measure.} 
Following R1's suggestion, we replaced total debt with long-term liabilities in the denominator of our gold exposure measure, $d_{j,t}$, now defined as the ratio of gold-clause bonds to long-term liabilities (corporate bonds, preferred shares, and bank loans) for firm $j$ at time $t$:
\begin{equation*}
d_{j,t} = \frac{\text{Gold-denominated long-term liabilities of firm $j$ in year $t$}}{\text{Long-term liabilities of firm $j$ in year $t$}}.
\end{equation*}

\item \textbf{Addressing endogeneity concerns.} Following R6's suggestion, we fix each firm's gold exposure at its 1930 value for all subsequent years to address potential endogeneity from strategic adjustments following Britain's departure from the gold standard in September 1931. Formally:
\begin{equation*}
\tilde{d}_{j,t} = \begin{dcases}
d_{j,t}, & \text{if } t \leq 1930 \\
d_{j,1930}, & \text{if } t > 1930
\end{dcases}.
\end{equation*}
This measure, $\tilde{d}$, is our main explanatory variable throughout the revised analysis.

\item \textbf{Empirical design.} Following R1's, R2's, and R6's suggestions, we now use the full 1926--1940 panel and estimate the effects of gold clause exposure and its reversal in a single regression.

\item \textbf{Debt overhang framing.} Following R1's and R2's suggestions, we have substantially revised the paper's conceptual framing. We no longer use the term ``leverage risk'' in the paper title or text and instead interpret the events of 1933--1934 through the lens of dynamic debt overhang theory \citep{Hennessy2004}. This framing more precisely captures the economic mechanism linking gold clause exposure to investment decisions.

\item \textbf{Firm-level uncertainty literature.} Following R2's suggestion, we have repositioned the paper within the firm-level uncertainty literature \citep{Bloom2009, Baker2016, Hassan2019}. The uncertainty surrounding the gold clause represents a policy shock of the type studied extensively in this growing literature. This reframing highlights that our gold clause exposure measure, $\tilde{d}_{j,t}$, captures firm-level variation in exposure to policy uncertainty---derived from balance sheet data rather than earnings call transcripts as in \citep{Hassan2019}. We are grateful to R2 for this suggestion, as it clarifies how our historical setting illustrates the real effects of policy uncertainty during a pivotal episode.
\end{itemize}
